\documentclass[a4paper,10pt]{article}
\newcommand\subdir{/home/koen/LaTeX-setup/}
\input{\subdir/LaTeX-files/imports.tex}
\input{\subdir/LaTeX-files/master-internship.tex}
\usepackage{\subdir/LaTeX-files/aastex}
\usepackage{natbib}
\bibliographystyle{\subdir/LaTeX-files/astroadstitle.bst}

%Set the title, Author and date
\title{Notes Week 30}
\author{Koen Essers}
\tutorial{Master Internship}
\date{\today}

%Set the header style
\header{\tutorialstring}

%Begin the document
\begin{document}
\setuplayout
\section{Checking differences between mixing methods}

Last week, I showed \emph{this} plot

\begin{figure}[ht]
	\centering
	\incplot{w29-large-grid-3}
\end{figure}

We would expect larger differences if we have a significant accretor core that is not mixed.
Let me investigate.

I quite quickly found the issue, which actually changes the results for every plot I showed last week that
was related to the binary abundance post-processing.

I accidentally took the final accretor mass \(M_\text{acc, f}\) as the total amount of accreted mass,
while this \emph{obviously} should be \(\Delta M_\text{acc} = M_\text{acc, f} - M_\text{acc, i}\).
This is an easy fix and now the plot above looks like this:

\begin{figure}[ht]
	\centering
	\incplot{w29-large-grid-3-fix}
\end{figure}

\section{Mean Molecular weight MESA vs Monash}

We saw that there is a difference between the mean molecular weight of the MESA
models and the initial mean molecular weight of the Monash models.

I want to understand where these differences come from.

Let me start by computing \(\mu \) from the \citet{asplund2009} solar abundances for the \emph{MESA} initial
Helium mass-fraction.

Then I will modify my \citet{asplund2009} abundances to the ``canonical'' helium mass-fraction as explained
in \citet{karakas2014}.

\begin{figure}[ht]
	\marginfignote{0}{The mean molecular weights of the \emph{MESA} models on the main sequence agree
		\emph{very} well with the mean molecular weights of the \citet{asplund2009} for both the
		standard \emph{MESA} initial helium mass-fraction as well as the ``canonical'' \citet{karakas2014}
		mass-fraction.
	}
	\marginfignote{12}{The \emph{MESA} models do not however, agree that well with the
		\citet{karakas2016} mean molecular weights.}
	\centering
	\incplot{w30-re}
\end{figure}

I plotted \emph{two} values for the initial TPAGB MESA model.
Below, I show some profiles during the
evolution of the \emph{MESA} model with the \citet{karakas2014} helium mass-fraction.

\begin{figure}[ht]
	\marginfignote{0}{There is a little dip just past the helium core boundary.
		The bottom blue dot in the
		plot above is sampled from this dip.
		The upper blue dot is sampled at \(R=1\), which is the value of
		most of the envelope, and probably the most relevant parameter.
	}
	\marginfignote{10}{The dip first forms towards the end of the GB already.
		This is \emph{after} the first dredge-up.
	}
	\centering
	\incplot{w30-mu-profiles}
\end{figure}

We can obviously also plot it with mass on the horizontal axis.
Then it becomes very
clear that the small dip is not the relevant value to use.

\begin{figure}[ht]
	\marginfignote{0}{The minimum is barely visible here.
		We can also clearly see the growth of the
		core between the GB and the TPAGB.
	}
	\centering
	\incplot{w30-mu-profiles-mass}
\end{figure}

I feel like the best thing to do would be to simulate everything with the \citet{karakas2014} initial helium
mass fraction.
This would however mean that I would need to simulate the single stars again, which
takes some time and manual help.

An easy first thing to do though is to simulate the accretors with the \citet{karakas2014} initial helium abundances.

This would at least ensure that the mean molecular weight profile of the main sequence stars should
align with the mean molecular weight of the Monash models.

However, the actual MESA binary models would still evolve with a different helium mass fraction.
I also think I should use the \citet{asplund2009} MS abundances for the accretor \emph{instead of}, what I do not,
use the initial envelope abundances of the TPAGB stars.

This is quite easy.

\begin{figure}[ht]
	\centering
	\incplot{w30-ms-vs-tpagb-ab}
\end{figure}

\nextpage

I simulated the main sequence stars with the canonical helium fractions from \citet{karakas2014}.

Here is the zoom on the mean molecular weight of the MESA helium mass fractions in gray,
and the new ones in color.

\begin{figure}[ht]
	\centering
	\incplot{w30-mu-core-2}
\end{figure}

The envelope has slightly larger mean molecular weight values.
This \emph{should} mean we get slightly lower mixing masses.

\begin{figure}[ht]
	\centering
	\incplot{w30-effects-of-mu-and-m_acc}
\end{figure}

This is indeed what we see.
It is much clearer for the second and third column if we zoom in the colors a bit.

\begin{figure}[ht]
	\centering
	\incplot{w30-effects-of-mu-and-m_acc-colorzoom}
\end{figure}

However, we also see that only certain regions are affected.

Below, I zoom in on a particular case of \(M_\text{acc}\) and \(\mu_\text{acc}\).
Lastly, I
\begin{figure}[ht]
	\marginfignote{0}{The diagonal lines show the approximate ages for the pre-MS and the MS.}
	\marginfignote{5}{I show the \(\mu \)-profiles of the cross formation of dots in the figure below.}
	\marginfignote{9}{The profiles clearly show where the differences between the two comes from.}
	\centering
	\incplot{w30-m_mix-grid-2}
\end{figure}

\begin{figure}[ht]
	\centering
	\incplot{w30-show-differences-mu-profiles}
\end{figure}
\nextpage

There are also these thin lines.
I zoom in one those below.
\begin{figure}[ht]
	\centering
	\incplot{w30-m_mix-grid-3}
\end{figure}

Again, the dots from a cross formation and the profiles are shown below.

\begin{figure}[ht]
	\centering
	\incplot{w30-show-differences-mu-profiles-2}
\end{figure}

I think I have the most important features in the difference plots.
I do not really know where the small bump
in the core composition comes from though.

\newsection{[X/Fe] abundances}

\begin{figure}[ht]
	\centering
	\incplot{w30-fe-ab}
\end{figure}

\begin{figure}[ht]
	\centering
	\incplot{w30-fe-ab-mass}
\end{figure}

\begin{figure}[ht]
	\centering
	\incplot{w30-fe-ab-mass-q}
\end{figure}

\begin{figure}[ht]
	\marginfignote{1}{All of these models have \([s/ \text{Fe}] \gg 0.25\) and would be classified
		as \(s\)-process enhanced star.}
	\marginfignote{6}{Clearly the lowest mass star has the lowest \([s/ \text{Fe}]\), which makes sense.}
	\centering
	\incplot{w30-grid-s-fe}
\end{figure}

\begin{figure}[ht]
	\centering
	\incplot{w30-hslss}
\end{figure}

\newsection{Broad radius grid}
I think it is interesting to look into the effects of dredge up on the final abundances of the accretor.
I want to simulate a MESA model for (nearly) every thermal pulse.
This should give me a good idea
of the effects of early RLOF interactions on the final abundances.

Below, I show the simple binary evolutions on which I based my MESA grid.
\begin{figure}[ht]
	\marginfignote{0}{Nearly every pulse should have a binary model.
		The last few models have multiple,
		because here the radius increases a lot.
	}
	\centering
	\incplot{w30-broad-R-grid-setup}
\end{figure}

\bibliography{master.bib}
\end{document}
